%----------------------------------------------------------------------------------------------------
% START OF CHAPTER
%----------------------------------------------------------------------------------------------------
\chapter{Feb 28}
\paragraph{Topics} 
\textit{ATPG. Prepositions. Essere, avere, and stare verbs. Present conditional verbs. 
Restaurant conversation}

\section{ATPG: All the Possibilities Game}

\begin{table}[ht]
	\centering
	\begin{tabular}{ r l  }
		\itb{prenotatzione} & \iti{reservation}               \\
		\hline
		the                 & \iti{la prenotatzione}          \\
		a                   & \iti{una prenotatzione}         \\
		two                 & \iti{due prenotatzione}         \\
		good                & \iti{una buona prenotatzione}   \\
		expensive           & \iti{una prenotatzione cara}    \\         
		new                 & \iti{una nuova prenotatzione}   \\
		big                 & \iti{una prenotatizone grande}  \\
		my                  & \iti{la mia prenotatzione}      \\
		your                & \iti{la tua prenotatzione}      \\
		his/her             & \iti{la sua prenotatzione}     
	\end{tabular}
	\caption{ATPG, Feb 28}
	\label{tab:tATPG-Feb28}
\end{table}

Irregular nouns (starting with \iti{s} + consonant or \iti{z}), see Table \ref{tab:tATPG-Feb28-part2}:

\begin{table}[ht]
\centering
    \begin{tabular}{ r l  | r l  | r l }
         \itb{scontrino} & \iti{receipt} & \itb{stato} & \iti{state} & \itb{zio} & \iti{uncle} \\
         \hline
         the        & \iti{lo scontrino}        & the       & \iti{lo stato}      & the   & \iti{lo zio} \\
         a          & \iti{uno scontrino}       & a         & \iti{uno stato}     & a     & \iti{uno zio} \\
         two        & \iti{due scontrini}       & two       & \iti{due stati}     & two   & \iti{due zii} \\
         good       & \iti{un buon scontrino}   & good      & \iti{un buon stato} & good  & \iti{un buon zio} \\
         expensive  & \iti{un scontrino caro}   & expensive & \iti{un stato caro} & 
	    dear\tablefootnote{\iti{Caro} means both `costly' and `beloved', like English `dear'.}  & \iti{un caro zio} \\
         new        & \iti{un nuovo scontrino}  & new       & \iti{una nuovo stato}  & new & \iti{un nuovo zio} \\
         big        & \iti{un scontrino grande} & big       & \iti{uno stato grande} & big & \iti{un zio grande} \\
         my         & \iti{il mio scontrino}    & my        & \iti{il mio stato}     & my  & 
	    \iti{mio zio\tablefootnote{Drop the article for family members: \iti{mio zio, mia madre}, not 
	    \iti{il mio zio, la mia madre}.}} \\
         your       & \iti{il tuo scontrino}    & your      & \iti{il tuo stato}  & your    & \iti{tuo zio} \\
         his/her    & \iti{il suo scontrino}    & his/her   & \iti{il suo stato}  & his/her & \iti{suo zio}
    \end{tabular}
    \caption{ATPG, Feb 28 - more unusual nouns}
    \label{tab:tATPG-Feb28-part2}
\end{table}

%--------------------------------------------------------------------------------
% Manual new page: ugly and not nice LaTeX, but makes the document flow better
% when rendered as PDF.
%--------------------------------------------------------------------------------
\newpage

\section{Grammar}
\subsection{Prepositions revisited}
Prepositions express spatial or temporal relations (in, under, on, between, etc.) or 
have a semantic roles (of, for). They precede a noun. See Table \ref{tab:tPrepositions-Revisited}.

\begin{table}[ht]
\centering
    \begin{tabular}{ r l  | r l   }
         with    & \iti{con} & for    & \iti{per}  \\ 
         between & \iti{tra} & under  & \iti{sotto} \\
         of      & \iti{di}\tablefootnote{A notable exception is \iti{di + Italia}, which 
	    becomes \iti{d'Italia}.}  & of the & \iti{del, della}  \\
         on      & \iti{su}  & on the & \iti{sul, sulla}  \\
        in the   & \iti{nel, nella} 
    \end{tabular}
    \caption{Common prepositions}
    \label{tab:tPrepositions-Revisited}
\end{table}

\newpage

\subsection{Prepositions with locations}
  Certain nouns, related to locations, require a specific preposition to represent `in'/`at'. 
	For example, Italian uses \iti{a}  for cities and towns, and \iti{in} for states and countries. 
	For example, \iti{Io vivo a Denver, in Stati Uniti}. More details in Table \ref{tab:tPrepositions-Locations}.

\begin{table}[ht]
\centering
    \begin{tabular}{ r l  | r l | r l }
           \textbf{in/at} home & \iti{a casa}   & \ldots the store & \iti{al negozio} & \ldots the restroom & \iti{in bagno} \\
           \ldots work & \iti{a lavoro} & \ldots the restaurant & \iti{al ristorante} & \ldots the church & \iti{in chiesa}\\
           \ldots Rome & \iti{a Roma}   & \ldots the mall & \iti{al centro commerciale} & \ldots Italy & \iti{in Italia}
    \end{tabular}
    \caption{Location prepositions - special cases}
    \label{tab:tPrepositions-Locations}
\end{table}
\paragraph{Examples}
\begin{itemize}
    \item I'm very happy because I am eating a pizza with Claudio at the restaurant. \\
    \iti{Io sono molto felice perché sto mangiando la pizza con Claudio al ristorante.}
    \item I'm at home because I am sick. \\
    \iti{Io sono a casa perché io sono malato.}
    \item I'm not at work because I am unwell. \\
    \iti{Io non sono a lavoro perché io sto male.}
    \item I'm buying lots of things at the shopping center. \\
    \iti{Io sto comprando molte cose al centro commerciale.}
    \item I'm driving Claudio's Maserati because I'm going to the store. \\
    \iti{Io sto guidando la Maserati di Claudio perché io sto andando al negozio.}
\end{itemize}


\subsection{Using the verbs \iti{stare} and \iti{essere}}
    The concept of \textbf{to be} spans two verbs in Italian: \iti{essere} and \iti{stare}. 
	For now, consider \iti{essere} as expressing presence, or a state that would use an 
	adjective, whereas \iti{stare} expresses a location, an ongoing action, or 
	a state that would use an adverb. 

    In colloquial US English, compare the response to the greeting `How are you?':
\begin{itemize}[noitemsep]
    \item I am happy. (Verb \iti{to be} + adjective \iti{happy}.)
    \item I am doing well. (Compound verb \iti{to be doing} + adverb \iti{well}.)
\end{itemize}

In Italian:
\begin{itemize}[noitemsep]
    \item I am happy. 
      (Verb: \iti{Io sono} + adjective \iti{felice}.)
    \item I am doing well. 
      (Compound verb: \iti{Io sto} + adverb \iti{bene}.)
\end{itemize}

In the US, people sometimes say ``I'm doing good'' but this is, strictly, ungrammatical 
English, and it wold be a mistake to carry that forward into Italian. 
It would be \uline{incorrect} to say this, because \iti{stare} takes an adverb not an adjective:
\begin{itemize}
    \item (Incorrect!) \iti{Io sto buono.} 
\end{itemize}

\paragraph{Examples}
\begin{itemize}
    \item How are you? \\
      \iti{Come stai?}
    \item I am fine. \\
      \iti{Io sto bene.}      
    \item I am not well. \\
      \iti{Io non sto bene.}
    \item I am unwell. \\
      \iti{Io sto male.}
    \item I am sick. \\
      \iti{Io sono malato.}
    \item I am happy. \\
      \iti{Io sono felice.}
    \item We are happy. \\
      \iti{Noi siamo felice.}      
\end{itemize}

\newpage

\subsection{The present progressive tense: \iti{stare} + gerund}
    The verb \iti{stare} is used to construct the \textbf{present progressive} tense, 
	where the speaker expresses something that is still happening. In English this 
	would be, for example, ``I am talking'' in contrast to ``I talk''. The `-ing' 
	form is a \textbf{gerund}. In Italian, gerunds end in \iti{-ando} for  \iti{-are} 
	verbs and  \iti{-endo} for \iti{-ire} and \iti{-ere} verbs .

\paragraph{Examples}
\begin{itemize}
    \item I am speaking. \\
      \iti{Io sto parlando.}
    \item She is eating. \\
      \iti{Lei sta mangiando.}      
    \item We are speaking. \\
      \iti{Noi stiamo parlando.}
    \item You are drinking. \\
      \iti{Tu stai bevendo.}
    \item Claudio is driving his car. \\
      \iti{Claudio sta guidando la sua macchina.}      
\end{itemize}

\newpage

\section{Vocabulary}
\subsection{Verbs}
\begin{table}[ht]
\centering
    \begin{tabular}{ r l | r l | r l }
         \itb{essere} & \textit{to be} & \itb{stare} & \textit{to be/to stay} & \itb{avere} & \textit{to have} \\
         \hline
         \itb{io}   & \iti{sono}   & \itb{io}   & \iti{sto}    & \itb{io}      & \iti{ho}      \\
         \itb{tu}   & \iti{sei}    & \itb{tu}   & \iti{stai}   & \itb{tu}      & \iti{hai}     \\
         \itb{lui/lei}  & \iti{è}  & \itb{lui/lei} & \iti{sta} & \itb{lui/lei} & \iti{ha}      \\
         \itb{noi}   & \iti{siamo} & \itb{noi}  & \iti{stiamo} & \itb{noi}     & \iti{abbiamo} \\
         \itb{voi}   & \iti{siete} & \itb{voi}  & \iti{state}  & \itb{voi}     & \iti{avete}   \\
         \itb{loro}  & \iti{sono} & \itb{loro} & \iti{stanno} & \itb{loro}    & \iti{hanno}   \\
    \end{tabular}
    \caption{\iti{essere, stare, avere}, Present tense}
    \label{tab:tVerbsPresentEssereStareAvere}
\end{table}

\paragraph{\iti{Mangiare} - to eat}
The verb \iti{mangiare} is pretty simple in the \textbf{present} tense 
(Table \ref{tab:tVerbsPresentAndPresentConditionalMangiare}). 
In the present \textbf{continuous} form (i.e., with \iti{stare}), the gerund 
\iti{mangiando} is used, e.g., \iti{io sto mangiando}, I am eating. 
The present \textbf{conditional} form, `would eat', is shown in the same table.
For example you might say, in a restaurant, 
\iti{Io mangerei la pasta}, literally, ``I would eat the pasta'', but more colloquially 
rendered as, ``I'll have the pasta''.

\begin{table}[ht]
\centering
    \begin{tabular}{ r l | l }
         \itb{mangiare} & \textit{to eat} &  \textit{could eat} \\
         \hline
         \itb{io}       & \iti{mangio}    & \iti{mangerei}      \\
         \itb{tu}       & \iti{mangi}     & \iti{mangeresti}    \\
         \itb{lui/lei}  & \iti{mangia}    & \iti{mangerebbe}    \\
         \itb{noi}      & \iti{mangiamo}  & \iti{mangeremmo}    \\
         \itb{voi}      & \iti{mangiate}  & \iti{mangereste}    \\
         \itb{loro}     & \iti{mangiano}  & \iti{mangerebbero} 
    \end{tabular}
    \caption{\iti{mangiare}, Present and present conditional tense}
    \label{tab:tVerbsPresentAndPresentConditionalMangiare}
\end{table}


\subsection{Words}
See Table \ref{tab:tVocabRestaurant}.

\begin{table}[ht] %[ht] % place table roughly [h]ere or else at the [t]op of page
\centering
    \begin{tabular}{ r l  | r l  | r l  }
          reservation & \iti{prenotatzione} &  table & \iti{tavolo} &  steak & \iti{bistecca} \\
          salad & \iti{insalata} & local wine & \iti{vino locale}  & bottle & \iti{bottiglia} \\
          sparkling & \iti{frizzante} &  dessert & \iti{dolce} &  bill & \iti{conto} 
    \end{tabular}
    \caption{Some restaurant terms}
    \label{tab:tVocabRestaurant}
\end{table}

\newpage

\subsection{Phrases}
\begin{itemize}
    \item Un tavolo per quattro persone, per favore.\\
      \iti{A table for four, please.}
    \item Ho una prenotatzione. \\
      \iti{I have a reservation.}
    \item Cosa mangiate? \\
      \iti{What would you like to eat?}
    \item Io proverei \ldots \\
      \iti{I'll try \ldots}
    \item Io mangerei \ldots questo piatto.\\
      \iti{I'll eat \ldots this dish. (E.g., pointing at something on the menu.)}
    \item Cosa da bere? \\
      \iti{Something to drink?}
    \item Tutto a posto? \\
      \iti{Everything OK?}
    \item Più pane, per favore. \\
      \iti{More bread, please.}
    \item Conto unico o separato? \\
      \iti{One check, or split?}
    \item Il piacere è mio. \\
      \iti{The pleasure is mine.}
\end{itemize}

\section{Conversation: \iti{Al ristorante}}
Myself, \iti{Io}, and a bunch of friends, \iti{i miei amici}, show up at a restaurant. 
We are met by \iti{il padrone} and are served by \iti{il cameriere}.
\begin{itemize}
    \item [Padrone] \iti{Buongiorno, tu hai una prenotatzione?}
    \item [Io] \iti{Sì, noi abbiamo una prenotatzione.}
    \item [Padrone] \iti{Perfetto.}
    \item [Cameriere] \iti{Cosa mangiate?}
    \item [Io] \iti{Io proverei le lasagne.}
    \item [amici] \iti{Un'insalata, per favore.}
    \item [amico] \iti{Io mangerei gli spaghetti carbonara con insalata.}
    \item [Io] \iti{La bistecca con patate.}
    \item [amico] \iti{Una pizza margherita, per favore.}
    \item [amica] \iti{Io proverei questo piatto, penne arrabbiata.}
    \item [amico] \iti{Io inizierei con un antipasto italiano.}
    \item [Cameriere] \iti{Perfetto! Grazie. Cosa da bere?}
    \item [Io] \iti{Due bottiglie di vino locale e due bottiglie di acqua frizzante.}
    \item [Padrone] \iti{Tutto a posto?}
    \item [Io] \iti{Sì, tutto a posto. Oh, più pane per favore.}
    \item [Cameriere] \iti{Caffè e dolce?}
    \item [Io] \iti{Caffè per favore.}
    \item [Io] \iti{Il conto, per favore}
    \item [Cameriere] \iti{Conto unico o separato?}
    \item [Io] \iti{Conto Unico. Io pago.}
    \item [$Tutti$] \iti{Grazie mille!}
    \item [Io] \iti{Il piacere è mio!}
    \item [Padrone] \iti{Grazie mille, e arrivederci.}
\end{itemize}

\section{Homework}
\begin{enumerate}
    \item I would eat, but I do not have time. \\
      \iti{Io mangerei, ma non ho tempo.}
    \item The wine in this restaurant is very good. \\
      \iti{Il vino in questo ristorante è molto buono.}
    \item I am at the store now. \\
      \iti{Io sono al negozio adesso.}
    \item I'm studying with Claudio's book. \\
      \iti{Io sto studiando con il libro di Claudio.}
    \item This table is for my friend. \\
      \iti{Questo tavolo è per il mio amico.}
    \item My name is Claudio. Pleased to meet you. \\
      \iti{Il mio nome è Claudio. Piacere.}
    \item I'm busy now, because I am at the bank. \\
      \iti{Sono occupato adesso, perché sono in banca.}
    \item More water, please. \\
      \iti{Più acqua, per favore.}
    \item My friend has family in Italy. \\
      \iti{La mia amica ha famiglia in Italia.}
    \item I am on vacation for two weeks. \\
      \iti{Io sono in vacanza per due settimane.}
\end{enumerate}
%----------------------------------------------------------------------------------------------------
% END OF CHAPTER
%----------------------------------------------------------------------------------------------------
